\fichatitulo{29 · REVISÃO}{Preprint v1 · 2024}{Embeddings na era dos LLMs}
{\small Nie et al. \textit{When Text Embedding Meets Large Language Model: A Comprehensive Survey}. Preprint v1 · 2024. \href{https://arxiv.org/pdf/2412.09165v1}{Fonte consultada}.\par}

\campo{Problema e objetivo}
Como LLMs e representações vetoriais de texto se relacionam?

\campo{Principal contribuição · síntese interpretativa}
Organiza o uso de LLMs para melhorar embeddings, produzir representações e interpretar informação codificada.

\campo{Método / natureza da evidência}
Revisão que articula geração e anotação de dados, modelos de representação e operações sobre embeddings.

\campo{Resultado ou argumento principal}
Evidencia diferentes papéis dos LLMs e desafios de treinamento, incluindo falsos negativos e contextos com poucos recursos.

\campo{Excertos-chave · seleção literal}
\begin{itemize}
\excerto{1}{LLMs as text embedders}{Taxonomia.}
\excerto{2}{false negatives}{Discussão de desafios.}
\end{itemize}

\campo{Limitações e análise crítica}
Análise da ficha: qualidade depende de tarefa, domínio e idioma. Sem avaliação local, uma taxonomia não justifica escolher um modelo de embeddings para discursos brasileiros.

\campo{Aproveitamento na dissertação · proposta}
Fundamentação: distinguir modelo gerador e modelo de representação; explicitar o critério experimental de seleção.

\registro{Fonte consultada nesta preparação: Resumo, taxonomia e trechos sobre desafios; consulta parcial. Chave: \texttt{nie2024textEmbeddingSurvey}.}
